% GENERATED FILE. Edit canonical lessons, then run npm run generate:books.
% canonical-source-hash: fnv1a64:3bb8a59b0236facf
% canonical-lessons: PT-C280-imprimir, PT-C280-publicacao, PT-C280-divulgar, PT-C280-edicao, PT-C280-artigo

\chapter{Nos jornais}
\label{ch:pt-nos-jornais}
\hlchaptermodality{280}

\begin{chapteropening}
\textbf{By the end of this chapter:} \emph{I can talk about printing a paper, publications, editions and articles, and making news public.}

Complete the last lesson of chapter 280: I can talk about printing a paper, publications, editions and articles, and making news public.
\end{chapteropening}

\section[imprimir]{imprimir — to print}

\label{lesson:PT-C280-imprimir}



\begin{quote}
Before the new one: say the Portuguese for the correspondent, then the Portuguese for the headline.
\end{quote}

\subsection*{You'll want to know: imprimir}
\textbf{imprimir} — \textquotedblleft{}to print\textquotedblright{}.

\begin{quote}
O jornal imprime cem mil exemplares por dia.
\end{quote}

\begin{quote}
\emph{The paper prints a hundred thousand copies a day.}
\end{quote}

Also on a home printer: \emph{imprimir um bilhete}. The past participle is \emph{impresso}.

\begin{cousinweb}
\emph{imprimir} is a learned borrowing from Latin \emph{imprimere} \textquotedblleft{}to press into, to stamp\textquotedblright{}, built on \emph{premere} \textquotedblleft{}to press\textquotedblright{} (Lewis \& Short).
\end{cousinweb}

\begin{grammarlens}[title={What you've built}]
One more word for telling, explaining and describing at B1.
\end{grammarlens}

\subsection*{Your turn}
\begin{itemize}
\raggedright
  \item \emph{Say it:} \emph{imprimir}
  \item \emph{Say it:} \emph{imprimir}, once more
  \item \emph{Say it:} a sentence of your own with \emph{imprimir}
  \item \emph{Recall:} say the Portuguese for the correspondent, then the Portuguese for the headline, then say \emph{imprimir} again
\end{itemize}

\begin{tcolorbox}[breakable,colback=teal!4,colframe=teal!35!black,title={Before you move on}]
What does imprimir mean? (\textquotedblleft{}To print\textquotedblright{}.) Say it once more.
\end{tcolorbox}
\section[a publicação]{a publicação — the publication}

\label{lesson:PT-C280-publicacao}



\begin{quote}
Before the new one: say the Portuguese for the headline, then the Portuguese for to print.
\end{quote}

\subsection*{You'll want to know: a publicação}
\textbf{a publicação} — \textquotedblleft{}the publication\textquotedblright{}.

\begin{quote}
A revista é uma publicação mensal.
\end{quote}

\begin{quote}
\emph{The magazine is a monthly publication.}
\end{quote}

Also a post on social media.

\begin{cousinweb}
\emph{publicação} is a learned borrowing from Latin \emph{pūblicātiō} \textquotedblleft{}a making public\textquotedblright{} (EtymDB), from \emph{pūblicāre}.
\end{cousinweb}

\begin{grammarlens}[title={What you've built}]
One more word for telling, explaining and describing at B1.
\end{grammarlens}

\subsection*{Your turn}
\begin{itemize}
\raggedright
  \item \emph{Say it:} \emph{a publicação}
  \item \emph{Say it:} \emph{a publicação}, once more
  \item \emph{Say it:} a sentence of your own with \emph{a publicação}
  \item \emph{Recall:} say the Portuguese for the headline, then the Portuguese for to print, then say \emph{a publicação} again
\end{itemize}

\begin{tcolorbox}[breakable,colback=teal!4,colframe=teal!35!black,title={Before you move on}]
What does a publicação mean? (\textquotedblleft{}The publication\textquotedblright{}.) Say it once more.
\end{tcolorbox}
\section[divulgar]{divulgar — to make public, to spread (news)}

\label{lesson:PT-C280-divulgar}



\begin{quote}
Before the new one: say the Portuguese for to print, then the Portuguese for the publication.
\end{quote}

\subsection*{You'll want to know: divulgar}
\textbf{divulgar} — \textquotedblleft{}to make public, to spread (news)\textquotedblright{}.

\begin{quote}
A empresa divulgou os resultados ontem.
\end{quote}

\begin{quote}
\emph{The company made the results public yesterday.}
\end{quote}

Formal; frequent on the news.

\begin{cousinweb}
\emph{divulgar} is a learned borrowing from Latin \emph{dīvulgāre} \textquotedblleft{}to spread among the people\textquotedblright{} (EtymDB), built on \emph{vulgus} \textquotedblleft{}the common people\textquotedblright{}.
\end{cousinweb}

\begin{grammarlens}[title={What you've built}]
One more word for telling, explaining and describing at B1.
\end{grammarlens}

\subsection*{Your turn}
\begin{itemize}
\raggedright
  \item \emph{Say it:} \emph{divulgar}
  \item \emph{Say it:} \emph{divulgar}, once more
  \item \emph{Say it:} a sentence of your own with \emph{divulgar}
  \item \emph{Recall:} say the Portuguese for to print, then the Portuguese for the publication, then say \emph{divulgar} again
\end{itemize}

\begin{tcolorbox}[breakable,colback=teal!4,colframe=teal!35!black,title={Before you move on}]
What does divulgar mean? (\textquotedblleft{}To make public, to spread (news)\textquotedblright{}.) Say it once more.
\end{tcolorbox}
\section[a edição]{a edição — the edition, the issue}

\label{lesson:PT-C280-edicao}



\begin{quote}
Before the new one: say the Portuguese for the publication, then the Portuguese for to make public.
\end{quote}

\subsection*{You'll want to know: a edição}
\textbf{a edição} — \textquotedblleft{}the edition, the issue\textquotedblright{}.

\begin{quote}
Comprei a edição de domingo.
\end{quote}

\begin{quote}
\emph{I bought the Sunday edition.}
\end{quote}

Also film or video editing.

\begin{cousinweb}
\emph{edição} is a learned borrowing from Latin \emph{ēditiō} \textquotedblleft{}a bringing forth, a publishing\textquotedblright{}, from \emph{ēdere} \textquotedblleft{}to put out\textquotedblright{}, built on \emph{dare} \textquotedblleft{}to give\textquotedblright{} (Lewis \& Short).
\end{cousinweb}

\begin{grammarlens}[title={What you've built}]
One more word for telling, explaining and describing at B1.
\end{grammarlens}

\subsection*{Your turn}
\begin{itemize}
\raggedright
  \item \emph{Say it:} \emph{a edição}
  \item \emph{Say it:} \emph{a edição}, once more
  \item \emph{Say it:} a sentence of your own with \emph{a edição}
  \item \emph{Recall:} say the Portuguese for the publication, then the Portuguese for to make public, then say \emph{a edição} again
\end{itemize}

\begin{tcolorbox}[breakable,colback=teal!4,colframe=teal!35!black,title={Before you move on}]
What does a edição mean? (\textquotedblleft{}The edition, the issue\textquotedblright{}.) Say it once more.
\end{tcolorbox}
\section[o artigo]{o artigo — the article}

\label{lesson:PT-C280-artigo}



\begin{quote}
Before the new one: say the Portuguese for to make public, then the Portuguese for the edition.
\end{quote}

\subsection*{You'll want to know: o artigo}
\textbf{o artigo} — \textquotedblleft{}the article\textquotedblright{}.

\begin{quote}
Li um artigo interessante sobre a saúde.
\end{quote}

\begin{quote}
\emph{I read an interesting article about health.}
\end{quote}

In grammar too: \emph{o} and \emph{a} are \emph{artigos}.

\begin{cousinweb}
\emph{artigo} is a learned borrowing from Latin \emph{articulus} \textquotedblleft{}a small joint, a part\textquotedblright{}, a little \emph{artus} \textquotedblleft{}a joint\textquotedblright{} (Lewis \& Short).
\end{cousinweb}

\begin{grammarlens}[title={What you've built}]
One more word for telling, explaining and describing at B1.
\end{grammarlens}

\subsection*{Your turn}
\begin{itemize}
\raggedright
  \item \emph{Say it:} \emph{o artigo}
  \item \emph{Say it:} \emph{o artigo}, once more
  \item \emph{Say it:} a sentence of your own with \emph{o artigo}
  \item \emph{Recall:} say the Portuguese for to make public, then the Portuguese for the edition, then say \emph{o artigo} again
\end{itemize}

\begin{tcolorbox}[breakable,colback=teal!4,colframe=teal!35!black,title={Before you move on}]
What does o artigo mean? (\textquotedblleft{}The article\textquotedblright{}.) Say it once more.
\end{tcolorbox}
